% -----------------------------------------------------------------------------
% Section 4: Security Analysis
% Paper 6 — The Living Constitution (Amendments)
% Security audit patch applied 2026-03-29:
%   Patch 6.1: Inherited Trust Assumptions subsection after security table
% -----------------------------------------------------------------------------

\section{Security Analysis}
\label{sec:security}

\subsection{Threat Model}

We consider an adversary who controls one or more of the three branches
and attempts to ratify an amendment that violates a Stone Clause
without obtaining the required supermajority. We also consider an
adversary who attempts to forge a ratification signature or replay a
previously ratified amendment.

\subsection{Security Properties}

\begin{table}[h]
\centering
\caption{Summary of Security Properties}
\label{tab:security}
\begin{tabular}{lll}
\hline
Property & Mechanism & Guarantee \\
\hline
Unforgeability & Ed25519 signatures & Stone Clauses cannot be silently overridden \\
Non-repudiation & Transparency log & Ratification is publicly verifiable \\
Replay resistance & Nonce + timestamp & Old amendments cannot be replayed \\
Capability disjointness & BTV type system & No branch can ratify alone \\
\hline
\end{tabular}
\end{table}

\subsection{Inherited Trust Assumptions}
\label{subsec:inherited-trust}

The Amendment Protocol inherits the trust assumptions of Papers~1--5.
Two inherited boundaries merit explicit acknowledgement in the
context of constitutional evolution:

\begin{enumerate}
\item \textbf{Key management uniformity.} The Tripartite
Ratification mechanism (Definition~\ref{def:tripartite}) assumes
that the signing keys of all three branches are protected to
HSM-grade standards. Paper~2~\cite{btv-t-repo} enforces this for
the Log Authority; Paper~1~\cite{btv2026} now reads the HMAC key
from an environment variable (\texttt{BTV\_HMAC\_KEY}), closing
the static-key limitation (L1) for production deployments.
An amendment that mandates HSM-backed key derivation for all
verdict seals---upgrading L1 from a deployment recommendation to a
constitutional requirement---is an example of
\emph{ordinary legislation} ($\Delta\mathcal{L}_{\mathsf{policy}}$,
Definition~\ref{def:amendment-classes}) that does not modify any
Stone Clause but materially strengthens the system's
non-repudiation posture.

\item \textbf{Demographic attribute integrity.} The Judicial
branch's ZK verification (Theorem~2 of
Paper~5~\cite{btv-cc}) assumes truthful demographic commitments
at ingestion time. Paper~3~\cite{btv-ar2026} proposes an Attested
Demographics extension to narrow this gap. Incorporating Attested
Demographics into the constitutional type schema---requiring
attribute-source signatures as a Stone Clause---would elevate
demographic integrity from an operational best practice to a
constitutional invariant, making its removal subject to Tripartite
Ratification.
\end{enumerate}

\noindent Neither boundary constitutes a vulnerability in the
Amendment Protocol itself; both are inherited from the layers the
protocol governs. We include them here for completeness, as any
security audit of the full series must evaluate the system as a
composed whole, not only at the amendment layer.
